\selectlanguage{portuguese}
\Language{pt}{Português}{Guia de utilização}
\firsttopic{Abrir, suspender e voltar}
O Inkline executa uma shell local no reMarkable 2. Use um Type Folio ou teclado
USB, trabalhe no tablet ou ligue-se por SSH ou Goblin Mosh. O texto e as imagens
são apresentados em tons de cinzento.

\begin{notice}
\key{Opt+RightAlt+T} abre o Inkline diretamente no tablet.\par
\key{Botão de alimentação:} prima e solte para suspender; prima novamente
para retomar. As shells e os buffers do editor ficam na RAM.
\end{notice}
A suspensão pausa o tablet sem fechar o Inkline nem guardar definições na
memória flash. As ligações de rede podem precisar de se restabelecer. O toque
que acorda o tablet não provoca outra suspensão. Após atualizar, feche e volte
a abrir o Inkline para ativar o novo comportamento do botão.

Para voltar aos cadernos, toque em \key{Quit} ou prima \key{Ctrl+Shift+Q} e
confirme. \code{exit} fecha um terminal; fechar o último devolve os cadernos.
A barra inferior contém Escape, histórico acima/abaixo, Quit e Settings.
\code{clear} apaga a breve ajuda inicial e o pequeno logótipo Goblin.

Se uma aplicação de atalho ocupar o ecrã, \key{Opt+RightAlt+T} regressa ao Inkline.
Numa emergência, mantenha \key{Opt+RightAlt+Backspace durante dois segundos}.
A recuperação fecha todos os terminais e perde o trabalho não guardado nessas
sessões. Após reiniciar o tablet, a interface inicial continua a ser a dos cadernos.

RightAlt é a tecla Alt/Opt à direita da barra de espaços. Mantenha-a premida com a tecla Opt separada e prima T, sem Ctrl. Super sozinho continua disponível para os seus atalhos.

\ProjectLinks

\topic{Teclas e posições de terminal}
O Type Folio tem uma tecla \key{Opt} separada entre Ctrl e Alt, e uma tecla
\key{Alt/Opt direita}. Têm funções diferentes.

\begin{keytable}
\key{Teclas} & \key{Ação} \\ \midrule
Opt separada + 1--0 & F1--F10 \\
Alt/Opt direita + 1--9 & Selecionar posição 1--9 \\
Alt/Opt direita + Left / Right & Posição anterior / seguinte \\
Alt/Opt direita + Tab & Escape \\
Alt/Opt direita + Up / Down & PageUp / PageDown \\
Alt/Opt direita + Backspace & Confirmar o fecho de todos os terminais \\
Qualquer Alt + Space & Teclado Unicode \\
Alt/Opt direita + C / V & Copiar seleção / colar \\
Ctrl+Shift+B & Ocultar / mostrar a barra inferior \\
\end{keytable}

No \key{Type Folio americano}, Alt/Opt direita+\key{0} escreve \code{+}.
Alt/Opt direita com a \key{tecla menos imediatamente à esquerda de Backspace}
escreve \code{=}. Com o método de entrada desligado, Shift+6 produz um
circunflexo literal: \code{c^2} mantém-se \code{c^2}. Não há atalhos de zoom.

Os teclados USB mantêm as teclas de função normais. O atalho da Opt separada
usa Meta quando o teclado USB disponibiliza essa tecla.

Existem até \key{nove posições independentes}. Selecionar uma vazia abre uma
shell. O indicador conta os terminais abertos: se apenas 1 e 9 estiverem
abertos, o 9 mostra \key{Terminal 2 / 2 · slot 9}. Fechar uma shell liberta a
posição sem alterar a numeração dos atalhos.

Cada terminal tem a sua shell, diretório, composição de entrada, imagens e
500 linhas de histórico. As definições de ecrã e de entrada são comuns.

\topic{Definições e ecrã}
Toque em \key{Settings}, o quinto botão inferior. Com a barra oculta, prima
\key{Alt+Space}, depois \key{F2}, ou toque em Settings no teclado Unicode.
Uma opção preenchida está ativa; o contorno tracejado indica o foco do teclado.
Tab move o foco; as setas e Enter alteram a opção.

\key{Caps Lock:} escolha Control (predefinição) ou Caps Lock.
\key{Tamanho:} faça pinça com dois dedos ou use os botões menos/mais.
Movimentos lentos ajustam um píxel de cada vez, de \key{6 a 48 píxeis}.

\key{Text darkness:} 50\% é a apresentação original. Valores superiores
escurecem os contornos. \key{Minimum contrast} separa a luminância do texto
e do fundo quando uma aplicação escolhe cores próximas; a predefinição é 35\%.

\begin{choices}
\key{Modo e-paper} & \key{Compromisso} \\ \midrule
Fast & Espera mínima; forma de onda de animação; mais imagens residuais \\
Balanced & Forma de onda de interface; agrupa saída em 32 ms \\
Crisp (predefinido) & Cinzentos mais limpos; forma de onda de conteúdo; mais lento \\
Mono & Preto e branco rápido; elimina os cinzentos das imagens \\
Saver & Agrupa em 120 ms; menos atualizações durante rajadas de saída \\
\end{choices}

Apenas as linhas alteradas são redesenhadas, mas o painel continua a precisar
de tempo. A poupança de bateria de Saver não foi medida. A atualização mantém
um modo guardado explicitamente.

\begin{notice}
Todas as definições ficam na \key{RAM} e são guardadas juntas ao sair normalmente
do Inkline. Levantar o dedo, fechar Settings ou suspender não escreve na flash.
Uma sessão sem alterações não escreve nada. Uma falha, encerramento forçado
ou perda de energia pode perder alterações ainda não guardadas.
\end{notice}

\topic{Unicode e métodos de entrada}
\key{Alt+Space} abre o teclado Unicode. Toque num caráter para o inserir;
o teclado permanece aberto. Toque em \key{Done} ou prima Escape para fechar.

As categorias incluem letras, marcas, números, pontuação, matemática/moedas,
símbolos, espaços, carateres de formatação e uso privado. Arraste na vertical
ou use setas, PageUp/PageDown, Home/End e a roda do rato.
\key{Tab} muda a categoria; Enter insere o caráter selecionado.

Escreva um valor Unicode hexadecimal, como \code{03bb}, seguido de Enter para
inserir λ. Backspace edita o valor. O círculo pontilhado das marcas combinantes
é apenas uma ajuda visual; só a marca é inserida. Espaços e carateres de
formatação têm etiquetas. O catálogo usa as tabelas Unicode do Qt, excluindo
controlos, valores não atribuídos, substitutos e não carateres. Alguns glifos
precisam de fontes não instaladas; o uso privado exige uma fonte adequada.

No seletor, prima \key{F2} ou toque em \key{Settings}, depois \key{Input method}:
\begin{choices}
Off — direct keyboard & Escrita americana normal \\
Romaji & Hiragana japonês a partir de entrada romanizada \\
Pinyin & Chinês através de pinyin \\
Zhuyin & Chinês com o mapa básico de zhuyin \\
Wubi 86 & Chinês através de códigos de forma \\
US-International & Acentos latinos comuns \\
\end{choices}

Escolha o resultado na barra de candidatos. Selecione
\key{Off — direct keyboard} no mesmo menu para voltar à entrada americana.
A escolha aplica-se a todos os terminais e é guardada ao sair do Inkline.

\topic{Toque, caneta e área de transferência}
\key{Dois dedos acima/abaixo:} percorrem o histórico. Arraste para baixo para
ver saída antiga e para cima para voltar ao prompt. Cada terminal mantém
500 linhas físicas em RAM, além do ecrã atual.

\key{Dois dedos à esquerda/direita:} mudam uma posição por gesto.
Levante os dedos antes de voltar a mudar. Afaste-os ou aproxime-os para
alterar o tamanho do texto.

\key{Um dedo:} arraste para selecionar texto. Alt/Opt direita+C copia;
Alt/Opt direita+V cola. Todos os terminais partilham uma área de transferência
em RAM até 128 KiB.

\key{Caneta:} as aplicações com relatórios de rato ativos recebem pressão,
movimento e libertação como eventos do botão esquerdo, usando o protocolo
negociado, incluindo SGR. Nos restantes casos, arrastar seleciona texto.
Mantenha \key{Shift} para forçar a seleção local numa aplicação que usa o rato.

\key{OSC 52} permite que programas do terminal, incluindo remotos por SSH,
leiam ou substituam esta área de transferência, separada da dos cadernos.
Alt/Opt direita+X copia e explica que a eliminação deve acontecer no editor:
OSC 52 não consegue apagar texto de um editor.

\key{Imagens:} os transportes kitty inline e por memória partilhada ficam
na RAM. Os modos de ficheiro e ficheiro temporário estão desativados.
Sixel funciona quando pedido explicitamente, mas não é anunciado automaticamente.
Os marcadores kitty, animações e alguns modos sixel continuam incompletos.

\code{clear} e a limpeza de todo o ecrã removem texto e imagens visíveis.
Imagens fora do ecrã são libertadas quando falta memória. Não há cache de imagens
em disco; imagens grandes continuam a consumir a RAM limitada do tablet.

\topic{Teclado USB e máquina de escrever}
Abra \key{Settings → USB (página 3)}. \key{Network} restaura Ethernet USB;
\key{Send keyboard} transforma o tablet num teclado para outro computador;
\key{Receive keyboard} recebe um teclado através de um adaptador OTG alimentado.
O Type Folio continua disponível. Instale o utilitário Python do site para
enviar teclas; não são necessários módulos Python adicionais.

Escolha Send keyboard, ligue USB-C, selecione o perfil e abra
\key{Open typewriter}. Coloque o foco num documento no computador e toque em
\key{Start}. Abre sempre em pausa. \key{Escape} ou \key{Stop} interrompe o envio
e liberta as teclas. Sair do ecrã, mudar de terminal, perder o foco, abrir uma
aplicação nativa ou suspender também pausa. Não há reenvio automático ao religar.

\begin{choices}
US / ASCII & Disposição US, Caps Lock desligado. Rejeita texto não ASCII. \\
Mac & Fonte de entrada Unicode Hex Input, Caps Lock desligado. Até U+FFFF;
rejeita caracteres dos planos suplementares. \\
Linux & Disposição US, Caps Lock desligado. A aplicação deve aceitar
Ctrl+Shift+U, código hexadecimal, Enter. Nem todas o permitem. \\
Windows & Instale WinCompose no PC, defina Compose como Right Alt e ative a
entrada Unicode. Disposição US, Caps Lock desligado. Envia Compose, u,
seis algarismos hexadecimais, Enter. \\
\end{choices}
WinCompose: \url{https://wincompose.info/}. As regras personalizadas não devem
substituir esta sequência. HID não oferece um canal universal de texto Unicode;
a configuração e os tipos de letra do destinatário importam. Teste um texto curto.

\key{Input method} oferece Off/US direto, Romaji, Pinyin, Zhuyin, Wubi 86 e
US-International. Só envia caracteres confirmados. Mudar de método pausa;
retome com Start. \key{Off} repõe a entrada US normal. \key{Unicode} ou
Alt+Space abre o seletor de caracteres para a saída USB.

Typewriter é um editor local cujo documento permanece na RAM durante o Inkline.
A escrita modifica o documento e, com o modo direto ativo, envia cada carácter
confirmado ao computador. Copiar, cortar, colar, selecionar e deslocar atuam no
editor local. \key{Files \& send} carrega e guarda UTF-8 com qualquer nome,
seleciona 5--80 caracteres por segundo, envia o documento completo e interrompe
o envio. Só Save ou a saída normal escrevem o documento. Sem nome, o ficheiro
oculto \path{~/.inkline-typewriter-draft} permite recuperar o texto após a saída.
\key{Exit typewriter} volta às definições USB.

\topic{Enviar ficheiros e restaurar a rede}
Os comandos \code{inkline-usb} e \code{keyboard-send} ficam em
\path{~/.local/bin}. \code{~/inkline usb}, \code{~/inkline type} e
\code{inkline-type} continuam como aliases compatíveis. As ferramentas Outpost ficam intactas.
\UsbExamples
Os ficheiros UTF-8 \key{não têm um limite fixo de tamanho}. São lidos em pequenos
blocos, sem guardar todo o documento na RAM nem criar ficheiros temporários.
Os ficheiros que permitem reposicionamento são verificados por inteiro antes
do envio e relidos; não altere a origem durante a operação. Os pipes são validados
à chegada, pelo que um erro pode surgir depois de texto já enviado. Remove-se
o BOM UTF-8 inicial e CRLF/CR passam a novas linhas. Enter e Tab atuam na aplicação ativa.

\code{--cps} (20 por omissão) ou \code{--delay-ms} define a velocidade; Unicode
precisa de mais relatórios. \code{--start-delay} permite mudar o foco. Dry-run
mostra relatórios sem USB. Ctrl+C ou \code{--stop} cancela e liberta as teclas.
Só funciona um emissor de cada vez. Uma fila interativa cheia pausa o envio.
Verifique o documento antes de repetir um envio parcial. Texto e estado temporário
não escrevem na memória flash. O perfil é guardado com as restantes definições
na saída normal.

\key{Network} ou \code{inkline-usb network} restaura Ethernet. Sair do Inkline,
incluindo após falha ou reinício de emergência, restaura a rede. O modo USB
não é guardado como preferência de arranque.
\UsbRecoveryExample
O temporizador é independente da shell. Mude de modo no tablet ou por SSH Wi-Fi,
nunca por SSH USB. Uma falha tenta recuperar Ethernet. Uma sessão PPP por satélite
ativa ou um gadget desconhecido mantém a porta; a regra também vale na saída.
Opt+RightAlt+Backspace reinicia o Inkline, mas fecha todos os terminais e perde
o trabalho não guardado.

\topic{Iniciar programas com atalhos}
Registe atalhos numa shell Inkline. Os argumentos são literais; invoque uma
shell explicitamente para usar a sua sintaxe. Os programas iniciados por atalho
leem \code{~/.bashrc} através de Bash.

Settings → Command shortcuts (página 2) permite atribuir um comando a uma letra. Escolha Terminal ou Native e-paper app e depois Save. Remove pede confirmação; Reset draft descarta as alterações. T está bloqueado. Só Save e a remoção confirmada escrevem no armazenamento. A ferramenta de linha de comandos gere as mesmas associações.

Todos os atalhos globais usam agora Opt+RightAlt, incluindo T e a recuperação mantendo Backspace premido. Use a tecla Opt separada do Folio; em teclados USB, use Windows/Command (Super). Ctrl+Alt continua disponível para o Emacs. Feche e reabra o Inkline após atualizar; a instalação preserva os terminais abertos.

\begin{shell}
~/inkline shortcut register e emacs
~/inkline shortcut register p python3
~/inkline shortcut list
~/inkline shortcut deregister p
\end{shell}

\key{Opt+RightAlt+E} abre o Emacs numa posição livre.
\code{~/inkline run PROGRAM ARG...} inicia sem registar um atalho.
Registar novamente substitui a associação; removê-la não afeta programas em
execução. As teclas seguem posições físicas americanas. Coloque pontuação
entre aspas e use caminhos completos para programas próprios, que devem
estar previamente instalados.

Para uma aplicação Qt/e-paper nativa que desenha o próprio ecrã:
\begin{shell}
~/inkline shortcut register a --epaper /home/root/my-app
\end{shell}
Uma aplicação nativa de cada vez pausa o Inkline e as suas shells. O lançador
fornece definições Qt e caches em RAM, descarta registos e retoma o Inkline
quando a aplicação termina. Use entrada Qt sem capturar exclusivamente o teclado.

\begin{notice}
\key{Opt+RightAlt+T é permanente.} Não pode ser registado nem removido.
Para a aplicação nativa e regressa aos terminais existentes. Não mantenha Ctrl premido para este atalho; a alteração de Caps Lock só se aplica dentro do Inkline.
\end{notice}

\key{Emergência:} mantenha Opt+RightAlt+Backspace dois segundos para parar a
aplicação e os seus descendentes e reiniciar o Inkline. Todos os terminais
fecham e perde-se trabalho não guardado. Aplicações root podem impedir a
recuperação capturando entrada, alterando serviços ou esgotando recursos.
Reiniciar o dispositivo é o último recurso.

\topic{Instalar, atualizar e guardar o guia}
O pacote destina-se a \key{reMarkable 2 com firmware 3.27} e usa as bibliotecas
Qt 6.8 e-paper do sistema. Transferências e verificação de integridade:

\InstallLink

O instalador verifica hardware, firmware, fontes, arranque do Qt, diretórios
temporários em RAM e swap desativado. As versões ficam em
\path{/home/root/.local/share/inkline}; o lançador de teclado inicia no arranque.
As atualizações mantêm os terminais abertos. \key{Feche e reabra o Inkline
para usar a nova versão.} As versões anteriores ficam para permitir regressar
e ocupam espaço adicional.

Este único PDF contém os nove idiomas e acompanha o pacote. A instalação tenta
importá-lo para \key{My files} pela interface web USB quando esta está ativa
e os cadernos estão abertos. Uma importação bem-sucedida é registada para não
duplicar o mesmo manual.

Se a importação automática não estiver disponível, ligue USB, feche o Inkline
e ative \key{USB web interface} nas definições de armazenamento dos cadernos.
Por SSH, execute:
\begin{shell}
~/inkline manual
\end{shell}
Também pode obter o PDF e importá-lo pelo navegador USB ou pela aplicação
reMarkable para computador/telemóvel. O importador nunca altera diretamente
a base dos cadernos nem interrompe terminais. O manual importado continua
em My files depois de desinstalar o Inkline.

Por SSH, volte aos cadernos ou remova o Inkline:
\begin{shell}
~/inkline stop
~/inkline uninstall
\end{shell}
A instalação de software opcional, incluindo o perfil LaTeX recomendado,
está documentada no site.

\topic{Funcionamento, armazenamento e RAM}
O Qt entrega teclas a um mapeamento comum para atalhos e Caps Lock.
O libghostty-vt codifica as teclas do terminal. Cada posição liga o programa
a um pseudoterminal (PTY) independente. As shells recebem
\code{HOME=/home/root}; o Bash interativo lê \code{~/.bashrc}.

A saída passa pelo analisador do terminal e pelo descodificador sixel.
O libghostty mantém células, cursor, ecrãs alternativos, histórico e imagens
kitty. O renderizador conserva uma imagem em cinzento e redesenha as linhas
alteradas; o backend e-paper do Qt atualiza o painel. Um daemon evdev independente
trata atalhos globais e alimentação; o systemd gere suspensão e retoma.

Cada terminal limita as alocações centrais a 64 MiB, as imagens kitty a
16 MiB por ecrã e as imagens sixel retidas a 16 MiB. Renderização e aplicações
consomem memória adicional. Posições vazias não alocam terminais.
Nove aplicações grandes podem exceder a RAM disponível.

Imagens temporárias, caches e sessões usam RAM. Software instalado, documentos,
definições guardadas ao sair e ficheiros escritos por aplicações usam armazenamento
persistente. O Inkline não impede os programas de escrever ficheiros.

O pacote Python opcional é Python padrão. Para reduzir escritas de cache,
adicione ao \code{~/.bashrc} do tablet:
\begin{shell}
export PYTHONDONTWRITEBYTECODE=1
export PIP_NO_CACHE_DIR=1
\end{shell}
Execute \code{source ~/.bashrc} ou abra outra shell. Use
\code{python3 -m pip install --no-compile PACKAGE} para evitar a compilação
de bytecode que o pip faz separadamente durante a instalação. O modo isolado
do Python ignora as suas variáveis de ambiente. O guia completo está no site.

O Inkline usa GPL-3.0-or-later. Os pacotes incluem arquivos de código-fonte
correspondentes e avisos das dependências. Os relatórios distinguem testes
automáticos de comportamentos físicos ainda por confirmar.
